% Generated by scripts/assemble.py; corrections belong in source/corrections.json.
\examhead{6}{北京交通大学 2025 年数学分析试卷}{13}
% p014-b01; source PDF page 14
\blockanchor{p014-b01}
\question{1}设 $f(x)$ 在 $\mathbb{R}$ 上可导, 若 $f(x)$ 有两个零点, 试证明在这两个零点之间必有 $f(x) + f'(x)$ 的零点.\par
% p014-b02; source PDF page 14
\blockanchor{p014-b02}
\question{2}设 $0 < x < 1$, 证明: $e^{2x} < \frac{x + 1}{1 - x}$.\par
% p014-b03; source PDF page 14
\blockanchor{p014-b03}
\question{3}已知 $a_0 = 1, a_{n+1} = \frac{1}{a_n^3 + 4}, n = 0, 1, 2, \cdots$. 证明:\par
% p014-b04; source PDF page 14
\blockanchor{p014-b04}
(1) $\{a_{n}\}$ 收敛.\par
% p014-b05; source PDF page 14
\blockanchor{p014-b05}
(2) $\{a_{n}\}$ 的极限是 $x^{4} + 4x - 1 = 0$ 的唯一正根.\par
% p014-b06; source PDF page 14
\blockanchor{p014-b06}
\question{4}求极限  $\lim_{n\to\infty}\left[\left(1+\frac{1}{n}\right)\left(1+\frac{2}{n}\right)\cdots\left(1+\frac{n}{n}\right)\right]^{\frac{1}{n}}$\par
% p014-b07; source PDF page 14
\blockanchor{p014-b07}
\question{5}已知  $u = f(r)$  二阶连续可微,  $r = \sqrt{x_{1}^{2} + x_{2}^{2} + \cdots + x_{n}^{2}}$ ,  $r \neq 0$ ,  $n \geq 1$ .\par
% p014-b08; source PDF page 14
\blockanchor{p014-b08}
(1) 试证明: $\sum_{k=1}^{n} \frac{\partial^2 u}{\partial x_k^2} = \frac{\mathrm{d}^2 u}{\mathrm{d}r^2} + \frac{n-1}{r} \frac{\mathrm{d}u}{\mathrm{d}r} (r \neq 0)$.\par
% p014-b09; source PDF page 14
\blockanchor{p014-b09}
(2) 求 $\sum_{k=1}^{n} \frac{\partial^2 u}{\partial x_k^2} = 0$ 的解.\par
% p014-b10; source PDF page 14
\blockanchor{p014-b10}
\question{6}求最大的负数 $A$ 和最小的正数 $B$, 使得不等式 $\frac{A}{xy} \leq \ln (x + y) \leq \frac{B}{x^2 + y^2}$ 对任意 $x > 0, y > 0$ 恒成立.\par
% p014-b11; source PDF page 14
\blockanchor{p014-b11}
\question{7}求重积分 $\iint_{D} |x - y + 2| \, \mathrm{d}x \, \mathrm{d}y$, 其中 $D: x^2 + y^2 \leq 4$.\par
% p014-b12; source PDF page 14
\par\Needspace{6\baselineskip}
\blockanchor{p014-b12}
\question{8}求曲面积分\par
\nopagebreak[4]
% p014-b13; source PDF page 14
\blockanchor{p014-b13}
\displayformula{\iiint_ {S} (x - y) \mathrm{d} y \mathrm{d} z + (y ^ {2} - z ^ {2}) \mathrm{d} z \mathrm{d} x + (z ^ {3} - x ^ {3}) \mathrm{d} x \mathrm{d} y.}
% p014-b14; source PDF page 14
\blockanchor{p014-b14}
其中 $S$ 是锥面 $\frac{x^2}{a^2} +\frac{y^2}{b^2} = z^2 (a,b > 0)$ 在 $z = 0$ 和 $z = 1$ 之间的部分，方向取下侧.\par
% p014-b15; source PDF page 14
\blockanchor{p014-b15}
\question{9}设 $f_0(x)$ 在 $[a, b]$ 上连续，记 $f_{n+1}(x) = \int_a^x f_n(t) \, \mathrm{d}t$，证明：函数项级数 $\sum_{n=0}^{\infty} f_n(x)$ 在 $[a, b]$ 上一致收敛.\par
% p014-b16; source PDF page 14
\blockanchor{p014-b16}
\question{10}计算级数 $\sum_{n=1}^{\infty} \frac{1}{2^n(2n-1)}$.\par
% p014-b17; source PDF page 14
\blockanchor{p014-b17}
\question{11}设对任意的  $n = 1, 2, \cdots, u_{n}(x)$  在 x = c 处连续，但级数  $\sum_{n=1}^{\infty} u_{n}(c)$  发散，证明：对任意的  $\delta > 0$ ，函数项级数  $\sum_{n=1}^{\infty} u_{n}(x)$  在  $(c, c + \delta)$  上非一致收敛.\par
% p014-b18; source PDF page 14
\blockanchor{p014-b18}
\question{12}设 $f(x)$ 有连续导数, $f_{n}(x) = n\left[f\left(x + \frac{1}{n}\right) - f(x)\right]$, 证明: $\{f_{n}(x)\}$ 在任意闭区间上均一致收敛.\par
